% ═══════════════════════════════════════════════════════════════════════════
% Master-thesis chapter OUTLINE: Methodology
% ---------------------------------------------------------------------------
% Your notes, brought into order. Conventions:
%   • Bullets = content to write in your own words.
%   • [KEEP] = your existing text snippet, already usable (polish wording).
%   • → Suggestion: = my answer to a question you asked in the notes.
%   • %-comments = organisational advice, not content.
% Figure placeholders compile via the same \figplaceholder macro as the
% results outline (defined there; redefine here if this file compiles alone).
% Processing/computer codes stay in their own chapter (your decision) — this
% chapter ends where the raw datasets exist on disk.
% ═══════════════════════════════════════════════════════════════════════════

\section{Methodology}
\label{sec:methods}

% Storyline of the chapter (one intro paragraph, then five subsections):
%   goal → equipment → cables (the "samples") → how the setup evolved
%   (two preliminary experiments, honestly reported as iterations) →
%   the final free-hanging-cable experiment that produced the main dataset.
% This order answers the classic methods questions in sequence: with what
% (equipment/materials), on what (cables), how did we get to the design
% (iterations), and what exactly was acquired (campaign + parameters).

% ─────────────────────────────────────────────────────────────────────────────
\subsection{Experimental Strategy and Overview}
\label{sec:methods:overview}

\begin{itemize}
  \item Goal restated in one sentence: measure the ground-to-cable strain
    transfer of free-hanging cable segments under controlled conditions,
    i.e.\ obtain the boundary displacement and the cable's own deformation
    field simultaneously → transfer function (cross-ref processing chapter).
  \item [KEEP] your opening idea-paragraph: initial plan = create a
    controlled reference strain field in a medium, attach free-hanging
    segments at two fixed points, measure strain below and on the fibre with
    the LDV, compare → transfer function.
  \item The approach evolved over three setups; state this openly as the
    design process (examiner-friendly): (1) aluminium beam — rigid medium,
    (2) ratchet strap — compliant medium, (3) direct boundary excitation of
    free-hanging cables (final; produced the primary dataset).
  \item One-sentence justification for the final design: instead of
    generating a strain field in a medium and coupling a cable to it, impose
    the boundary displacement directly on the cable — removes the medium as
    an uncontrolled intermediate step and makes $\delta L$ a measured input.
\end{itemize}

% Roadmap table — lets you compress the two failed setups drastically while
% still "showing what was done and what data was acquired".
\begin{table}[H]
  \centering
  \begin{tabularx}{\textwidth}{l l l X}
    \toprule
    \textbf{Setup} & \textbf{Medium / specimen} & \textbf{Source signals} &
    \textbf{Outcome} \\
    \midrule
    1. Aluminium beam & beam on foam, cable planned on top &
      12\,Hz Ricker &
      beam moved quasi-rigidly; no resolvable differential strain
      ($\lambda \gg$ beam) → abandoned \\
    2. Ratchet strap & tensioned 1.95\,m strap, cable segments on spacers &
      12\,Hz Ricker; 1–500\,Hz linear sweep (9\,s) &
      strap-scale velocities characterised (standing-wave analysis);
      small-scale strain inconsistent → abandoned \\
    3. Free-hanging cables & cable clamped shaker $\leftrightarrow$ fixed end &
      1–500\,Hz log sweep (10\,s); mono-frequency ramps &
      primary dataset (98 recordings — verify count, see
      §\ref{sec:methods:campaign}) \\
    \bottomrule
  \end{tabularx}
  \caption[Evolution of the experimental setups]{Evolution of the
    experimental setups. % fill/adjust outcomes in your words
  }
  \label{tab:methods:evolution}
\end{table}

% ─────────────────────────────────────────────────────────────────────────────
\subsection{Laboratory Environment and Equipment}
\label{sec:methods:equipment}

\subsubsection{Scanning Laser Doppler Vibrometer (MATRIX / CIWE)}
\begin{itemize}
  \item [KEEP] your paragraph: CIWE's MATRIX research platform
    (\textit{Machine for Time Reversal and Immersive experimentation}),
    robot-mounted three-component scanning LDV by Polytec; measures particle
    velocity at designated points; repeatable source → spatially dense
    acquisition ("virtual array", one point at a time).
    % → Suggestion for the missing citation: cite the CIWE/MATRIX platform
    %   paper if one exists (ask Johan/Simone — there is an ETH CIWE
    %   description paper) + the Polytec 3-D scanning-vibrometer datasheet
    %   for instrument specs (also covers the 50 cm standoff below).
  \item Working principle one-liner (Doppler shift of backscattered laser
    light → line-of-sight velocity; three heads → 3-D velocity vector) +
    citation to Polytec documentation.
  \item Optimal focal standoff $\approx$ 50\,cm from the target [KEEP —
    source: Polytec datasheet, see suggestion above].
  \item Surface preparation for reflectivity: reflective tape on the
    aluminium beam / strap; \textbf{chalk spray} on the cables of the main
    experiment [KEEP].
  \item Stacking: each point measured repeatedly per source repetition,
    \textbf{median} stack (robust against dropouts) — state the number of
    averages here or in the acquisition table (5).
  \item \textbf{Coordinate systems} (your note, ordered):
    \begin{itemize}
      \item robot arm has a global coordinate system;
      \item per experiment, four reference points define a local system —
        aluminium-beam experiment: the four top corners of the beam;
        all later experiments: the four top corners of the aluminium
        profile frame;
      \item all reported X, Y, Z coordinates are in this local system;
      \item known offset between local system and the shaker-head origin,
        corrected in preprocessing: $x$: 15.9\,cm, $y$: 6\,cm, $z$: 6\,cm.
        % → Suggestion to your 14.9-vs-15.9 question: the analysis code
        %   (ldv_analysis/io.py) subtracts 0.159 m — every processed result
        %   uses 15.9 cm. Recommend: report the tape-measured value with its
        %   uncertainty, state that the applied offset (15.9 cm) was
        %   verified against the data (e.g. position of the first sensor /
        %   shaker connector in the scans), and attribute the 1 cm
        %   difference to the connector geometry. Check the physical
        %   measurement once more before writing.
    \end{itemize}
  \item Point-grid definition [KEEP]: first and last point marked manually in
    the LDV software; intermediate points interpolated by successive midpoint
    insertion → equidistant spacing.
\end{itemize}

\begin{figure}[htb]
  \centering
  \figplaceholder{Sketch: global robot frame → 4 local reference points on
    the alu-profile corners → local XYZ axes; shaker-head origin with the
    (15.9, 6, 6) cm offset arrows; laser heads + 50 cm standoff}
  \caption{Coordinate-system definition of the LDV acquisition. [Your
    sketch note from the LDV subsection.]}
  \label{fig:methods:coords}
\end{figure}

\subsubsection{Shaker Source and Signal Chain}
\begin{itemize}
  \item [KEEP] Data Physics GW-V4/PA30E shaker, DC–14\,000 Hz; driven by the
    PA30E amplifier; trigger/signal from the LDV control computer
    (synchronised acquisition).
  \item Signal chain in one sentence: waveform defined digitally (voltage) →
    amplifier (gain setting R2) → shaker head displacement; the \emph{actual}
    head motion is measured by the LDV on the shaker connector (the
    \texttt{\_SHAKER} reference recordings) — so no assumption about the
    shaker's own transfer function is needed.
    % → Suggestion for your "displacement vs velocity" note: don't discuss
    %   the shaker's internal displacement response in the text — it is
    %   irrelevant precisely because the reference channel measures the
    %   realized velocity, which the processing integrates to displacement.
    %   One sentence as above suffices.
  \item Frequency band of interest: seismology-motivated, primarily
    $<$ 50 Hz quasi-static regime, but sweeps to 500 Hz to capture the
    resonance-governed dynamic regime [KEEP, tie to theory chapter].
\end{itemize}

% Source-signal table — replaces the loose paragraphs; the "purpose/used in"
% column tells the story of the iterations without extra prose.
\begin{table}[H]
  \centering
  \begin{tabularx}{\textwidth}{l l X}
    \toprule
    \textbf{Signal} & \textbf{Parameters} & \textbf{Purpose / used in} \\
    \midrule
    Ricker wavelet & $f_c = 12$\,Hz &
      transient test signal; alu-beam and strap setups \\
    Linear sweep & 1–500\,Hz, 9\,s &
      strap experiment; standing-wave characterisation \\
    Logarithmic sweep & 1–500\,Hz, 10\,s sweep, 12.5\,s record incl.\
      buffer &
      \textbf{main dataset}; equal number of cycles per octave → low
      frequencies get full periods and sweep slowly [KEEP your reasoning] \\
    Mono-frequency ramps & 8 tones (5–400\,Hz), 3.5\,s each, linear
      amplitude ramp $0\!\to\!1$ + 0.5\,s fade, 38\,s record &
      linearity test (amplitude dependence of the coupling) \\
    \bottomrule
  \end{tabularx}
  \caption[Source signals]{Source signals used across the experimental
    campaign.}
  \label{tab:methods:signals}
\end{table}

\begin{figure}[htb]
  \centering
  \figplaceholder{(a) programmed log-sweep time series; (b) measured
    spectrogram of the shaker reference channel (shows realized sweep +
    harmonics)}
  \caption{Main excitation signal: programmed waveform and measured
    response. [Your "plots of programmed sweep + spectrograms" note.]}
  \label{fig:methods:sweep}
\end{figure}

% ─────────────────────────────────────────────────────────────────────────────
\subsection{Fibre-Optic Cables Under Test}
\label{sec:methods:cables}
% Your instinct in the notes was right: present the cables (incl. measured
% properties) BEFORE the experiments — they are the "materials".

\subsubsection{Cable Selection and Description}
\begin{itemize}
  \item Seven cables spanning a wide range of diameter (0.9–5.5 mm),
    stiffness and construction — chosen to span several decades of the
    coupling parameter $\Theta$ (cross-ref theory chapter).
  \item Per cable 1–2 sentences: construction (jacket material, armouring,
    number of fibres, loose-tube vs tight-buffered) — from spec sheets where
    available [your "give some descriptions" note; the spec-sheet names in
    Table~\ref{tab:methods:cabletypes} contain most of this].
  \item State the renumbering explicitly: cables are labelled C1–C7 sorted by
    descending diameter; the raw data files and analysis code use the
    original numbering (given in parentheses).
    % → Important consistency note: the dataset names and ldv_analysis
    %   config use the ORIGINAL numbering (e.g. code "Cable5" = thesis C1).
    %   Keep the mapping column in every table, or you (and readers of your
    %   data) will mislabel cables. Consider a footnote in the results
    %   chapter too.
\end{itemize}

% [KEEP] your cable-types table (unchanged, one suggested extra column):
\begin{table}[H]
  \centering
  \begin{tabularx}{\textwidth}{l l l X}
    \toprule
    \textbf{Cable ID} & \textbf{Name} & \textbf{Description} &
    \textbf{Construction} \\
    \midrule
    C1 (formerly C5) & Silixa PS-2S2M-1PU055-01 & Black, Silixa Carly &
      \textit{[jacket, fibres, buffering]} \\
    C2 (C1) & Metal-armored, unknown manufacturer & Black, Metal Core &
      \textit{[\ldots]} \\
    C3 (C2) & Silixa PS-48-2PE038-01-B & Black, Silixa BRUsens &
      \textit{[\ldots]} \\
    C4 (C3) & FIBRASENS 2OS2-FS02-19/31-T-BU & Blue, Rectangular &
      \textit{[\ldots]} \\
    C5 (C6) & Plus Corning / TLC 3\,mm OD OFNR SMF-28 Ultra & Yellow, Thick
      Carly & \textit{[\ldots]} \\
    C6 (C4) & HUBER+SUHNER 4$\times$9/125 LowBend G657.A2 LSFH 2.3\,mm &
      Yellow, Mid Dübendorf & \textit{[\ldots]} \\
    C7 (C7) & ThorLabs SMF-28-100, \diameter 900\,µm Hytrel jacket & Yellow,
      Thin Carly & \textit{[\ldots]} \\
    \bottomrule
  \end{tabularx}
  \caption[Cable types]{Tested fibre-optic cables, sorted by diameter.
    Original labels used in the raw data in parentheses.}
  \label{tab:methods:cabletypes}
\end{table}

\begin{figure}[htb]
  \centering
  \figplaceholder{Photo of the seven cables side by side (with scale bar /
    ruler), labels C1–C7}
  \caption{The seven tested cables. [Your photo note.]}
  \label{fig:methods:cablephoto}
\end{figure}

\subsubsection{Geometric and Mass Properties}
\begin{itemize}
  \item [KEEP] dimensions from spec sheets, coil markings, or caliper
    measurement (source per cable is the "Diameter" column note in
    Table~\ref{tab:methods:cableparams}).
  \item Cross-section: $A = \pi r^2$ (circular); C4 rectangular
    $A = w \cdot h$ with $w = 3.1$\,mm, $h = 1.9$\,mm.
  \item [KEEP] bulk-density procedure: sample piece cut (30\,cm; 2\,m for the
    thin C7), length measured (±2\,mm estimated), weighed on a METTLER
    AE~200 analytical scale (reading accuracy 0.1\,mg; your estimated
    uncertainty ±0.5\,mg).
  \item Equation to state (your note asked for it):
    \[ \rho \;=\; \frac{m}{V} \;=\; \frac{m}{L \cdot A}
       \qquad\text{(equivalently } \rho = w_\ell / A
       \text{ with linear weight } w_\ell = m/L\text{)} \]
    and straightforward relative-uncertainty propagation
    \[ \frac{\sigma_\rho}{\rho} \;=\;
       \sqrt{\left(\frac{\sigma_m}{m}\right)^{\!2}
           + \left(\frac{\sigma_L}{L}\right)^{\!2}
           + \left(\frac{\sigma_A}{A}\right)^{\!2}} , \]
    where $\sigma_A$ follows from the diameter uncertainty
    ($\sigma_A/A = 2\,\sigma_d/d$). With $\sigma_L = 2$\,mm on 30\,cm the
    length term ($\approx$0.7\,\%) dominates over the mass term
    ($\approx$0.01\,\%) — the diameter term likely dominates overall for the
    caliper-measured cables.
\end{itemize}

% Weighing-table template (your spec: ID, length, mass, g/m, density) —
% masses filled from your notes; g/m computed as m/L; VERIFY before use.
\begin{table}[H]
  \centering
  \begin{tabularx}{\textwidth}{l l l l l X}
    \toprule
    \textbf{ID} & \textbf{Sample length} & \textbf{Mass / g} &
    \textbf{Linear weight / g\,m$^{-1}$} & \textbf{Area / mm$^2$} &
    \textbf{Bulk density / kg\,m$^{-3}$} \\
    \midrule
    C1 & 30\,cm & 7.4928 & 24.98 & \textit{[23.76]} & 1056.26 \\
    C2 & 30\,cm & 7.7080 & 25.69 & \textit{[11.34]} & 2263.14 \\
    C3 & 30\,cm & 3.9266 & 13.09 & \textit{[11.34]} & 1146.27 \\
    C4 & 30\,cm & 2.7131 & 9.04  & \textit{[5.89]}  & 1528.01 \\
    C5 & 30\,cm & 1.9062 & 6.35  & \textit{[7.07]}  & 943.13 \\
    C6 & 30\,cm & 1.4862 & 4.95  & \textit{[4.15]}  & 1203.44 \\
    C7 & 2\,m   & 1.5773 & 0.79  & \textit{[0.64]}  & 1257.52 \\
    \bottomrule
  \end{tabularx}
  \caption[Bulk-density determination]{Bulk-density determination
    (uncertainties: length ±2\,mm, mass ±0.5\,mg, both estimated). Full
    weighing sheets in Appendix~X. % [your appendix note]
    }
  \label{tab:methods:density}
\end{table}
% → CHECK before writing: your parameter-table note lists C5 at 6.67 g/m,
%   but 1.9062 g / 0.30 m = 6.35 g/m. One of the two numbers (mass or g/m)
%   is stale — recompute C5's density from whichever is the raw measurement.
%   All other cables are consistent (24.98≈25, 25.69≈25.67, 13.09≈13,
%   9.04≈9, 4.95≈5, 0.79≈0.8).

\subsubsection{Elastic Modulus}
\begin{itemize}
  \item [KEEP] your full tensile-test paragraph — it is already good methods
    text: Zwick Roell Z010 at the ETH Institute for Building Materials;
    three 15\,cm samples per cable; clamping jaws 9\,cm apart; 5\,N preforce
    to straighten; Zwick Roell extensometer (WN 180765, 50\,mm baseline,
    roll–blade arrangement) measuring elongation → force–deformation curve.
  \item [KEEP] evaluation: $\sigma = F/A$, $\varepsilon = \Delta L / L_0$,
    $E = \sigma/\varepsilon$ from the linear regime; linear regression
    between the points at 40\,\% and 80\,\% of peak force; mean of three
    samples reported.
  \item [KEEP] the engineering-stress remark (constant initial
    cross-section assumed, no necking correction) — one honest sentence.
  \item Report $E$ as mean ± std of the three samples (the std propagates
    into $\Theta_\mathrm{mat}$ — cross-ref uncertainty section of the
    processing chapter).
  \item Appendix: per-sample values + one representative force–elongation
    curve per cable [KEEP your appendix plan].
\end{itemize}

\begin{figure}[htb]
  \centering
  \figplaceholder{Photo/sketch of the tensile tester with extensometer
    attached to a cable sample}
  \caption{Tensile-test setup (Zwick Roell Z010 with extensometer).}
  \label{fig:methods:tensile}
\end{figure}

\begin{figure}[htb]
  \centering
  \figplaceholder{One representative cable: three force–elongation curves +
    the 40–80\,\% regression fits (other cables → appendix)}
  \caption{Force–elongation curves and linear fits used for the elastic
    modulus. [Your figure note.]}
  \label{fig:methods:ecurves}
\end{figure}

% Consolidated parameter table — your CABLE PARAMETERS table, completed with
% the E column. E values below are the means (± std) of the three tests as
% stored in ldv_analysis/config.py — VERIFY against your lab sheets
% (remember: config.py uses the ORIGINAL cable numbering).
\begin{table}[H]
  \centering
  \begin{tabularx}{\textwidth}{l l l l X}
    \toprule
    \textbf{ID} & \textbf{Diameter / mm} & \textbf{Weight / g\,m$^{-1}$} &
    \textbf{Bulk density / kg\,m$^{-3}$} & \textbf{Young's modulus / GPa} \\
    \midrule
    C1 (C5) & 5.5 (spec sheet) & 25 & 1056.26 & 8.87 $\pm$ 0.63 \\
    C2 (C1) & 3.8 (caliper) & 25.67 & 2263.14 & 28.62 $\pm$ 2.90 \\
    C3 (C2) & 3.8 (on cable) & 13 & 1146.27 & 11.13 $\pm$ 0.90 \\
    C4 (C3) & 1.9 $\times$ 3.1 (from coil) & 9 & 1528.01 & 4.21 $\pm$ 0.19 \\
    C5 (C6) & 3 (spec sheet) & 6.67 & 943.13 & 2.79 $\pm$ 0.18 \\
    C6 (C4) & 2.3 (spec sheet) & 5 & 1203.44 & 6.79 $\pm$ 0.39 \\
    C7 (C7) & 0.9 (spec sheet) & 0.8 & 1257.52 & 2.06 $\pm$ 0.16 \\
    \bottomrule
  \end{tabularx}
  \caption[Cable properties]{Cable properties. Diameters from spec sheets,
    coil markings or caliper as indicated; density from
    Table~\ref{tab:methods:density}; $E$ = mean ± std of three tensile
    tests.}
  \label{tab:methods:cableparams}
\end{table}

% → Suggestion to your "eigenfrequency — maybe not relevant" note: drop the
%   f₁ formula from this chapter. It already lives in the theory/processing
%   chapters (clamped–clamped Euler–Bernoulli, used by the modal analysis),
%   and the prediction-vs-observation comparison is a RESULT (modal section
%   of the results outline). Here, at most one sentence: "the measured
%   properties also feed the theoretical resonance estimates (Section …)".

% ─────────────────────────────────────────────────────────────────────────────
\subsection{Preliminary Experiments}
\label{sec:methods:preliminary}
% → Suggestion to your "how do I show the failed ideas" question: give each
%   preliminary setup the same fixed mini-structure — (a) setup in ≤1
%   paragraph, (b) what was observed (1 figure max in the main text),
%   (c) why it was insufficient, (d) what the final design inherits from it.
%   That framing turns "failures" into design iterations; move acquisition
%   tables and extra figures to the appendix.

% Common frame first (used by all setups):
\begin{itemize}
  \item [KEEP] general frame: aluminium profile (Kanya, 120 $\times$
    2200\,mm) mounted on two wooden trestles with foam in between for
    mechanical decoupling; trestles connected at the bottom by a wooden beam
    for stability/weight; shaker bolted to one end of the profile, head
    horizontal, aligned with a spirit level.
\end{itemize}

\begin{figure}[htb]
  \centering
  \figplaceholder{Photo of the general frame: profile on trestles, shaker at
    the end (annotated)}
  \caption{Common experimental frame.}
  \label{fig:methods:frame}
\end{figure}

\subsubsection{Aluminium Beam}
\begin{itemize}
  \item [KEEP] setup: beam on foam on the rail (elevation + decoupling);
    3-D-printed connection cylinder bolted to the shaker and glued to the
    beam end with Crystalbond adhesive (Agar Scientific AGB7313, ID
    555-HMP); reflective tape at 2\,cm spacing.
  \item [KEEP] acquisition: 12\,Hz Ricker; fs = 100\,kHz; 1\,cm spatial
    sampling → 199 points; LDV field of view limited → beam scanned from six
    robot positions, datasets stitched.
  \item [KEEP] physics of the failure: P-wave velocity $\approx$ 5000\,m/s
    (cite your research-module report) → wavelength $>$ 400\,m $\gg$ 2\,m
    beam → the beam translates quasi-rigidly; differential displacement
    along the beam (the strain) is below what is resolvable.
  \item Show it, don't argue it: one figure of measured axial displacement /
    strain along the beam ("the beam was pushed, not strained") [your plot
    note].
  \item Lesson carried forward: a rigid medium cannot produce a measurable
    reference strain field at these frequencies and lengths → need either a
    compliant medium (→ strap) or direct boundary excitation (→ final).
\end{itemize}

\begin{figure}[htb]
  \centering
  \figplaceholder{Alu-beam result: velocity/displacement along the beam
    showing near-uniform motion (no differential strain)}
  \caption{Aluminium-beam test: the beam responds quasi-rigidly. [Supports
    the abandonment argument.]}
  \label{fig:methods:alubeam}
\end{figure}

\subsubsection{Ratchet Strap}
\begin{itemize}
  \item [KEEP] setup: 2\,cm-wide ratchet strap ($L = 1.95$\,m) on a wooden
    beam (screw-clamped to the profile) for height alignment with the
    shaker; strap tensioned, far end clamped at $\approx$190\,cm
    (screw clamp to wood + profile); cardboard strip with packing tape under
    the strap to reduce friction; strap glued to the shaker head
    (superglue); reflective tape every 10\,cm (2–192\,cm).
  \item [KEEP] cable mounting for the transfer test: two flat spacers
    ($\approx$2\,mm high, 2$\times$2\,cm strap pieces on double-sided tape)
    at defined positions (35–40\,cm and 95–100\,cm; a third position can be
    dropped — your note); 9\,cm cable segments taped (electrical tape)
    across them → free-hanging gap of 5 or 12\,cm; LDV lines on the cable
    and directly next to it on the strap (reference), 0.5\,cm spacing,
    ±2\,cm beyond the gap.
  \item [KEEP] acquisitions: cables C1–C4 (note: original numbering? verify),
    12\,Hz Ricker + 1–500\,Hz linear sweep (9\,s); high spatial sampling
    (fill exact value); acquisition-parameter table → appendix [your note].
  \item \textbf{Strap characterisation} (condense your standing-wave report
    to $\approx$½ page + 1–2 figures; full report → appendix; it is your
    own unpublished work — no plagiarism issue, but keep the derivation
    cross-referenced):
    \begin{itemize}
      \item observation: low-amplitude zones migrating source-ward with
        frequency during the sweep = nodes of a standing wave (fixed-end
        reflection, superposition $u = -2A\cos(kx)\sin(\omega t)$);
      \item node positions ↔ frequency give the phase velocity:
        $c(f) = 2f\,(L - x_1)$ for the single interior node;
      \item node tracking via Hilbert-envelope minima (amplitude-sensitive,
        not robust) and via the spatial gradient of the unwrapped
        instantaneous phase (the $\pi$ polarity flip across a node —
        robust); 2nd-order fit of node trajectories;
      \item result: $c \approx 950$\,m/s around 250\,Hz rising to
        $\approx$1150\,m/s at 500\,Hz (weakly dispersive); confirmed
        independently by the f–k spectrum slope ($\approx$1000\,m/s);
      \item purpose in the thesis: proves the strap carries a usable strain
        field at metre scale and documents the wave regime the cable
        segments were exposed to.
    \end{itemize}
  \item [KEEP] why abandoned: amplitudes not reproducible between
    acquisitions; suspected small-scale material heterogeneity (used strap),
    non-constant source output, and insufficient small-scale strain
    resolution — strain from spatial differentiation amplified high
    wavenumbers (striping artefacts); wavenumber muting acted only as
    spatial smoothing. Despite many processing attempts, small-scale strain
    was not reliable → new approach.
    % → Suggestion for your open TODO ("check if the new M1 gradient method
    %   fixes the noise"): keep it out of the methods text. If you do run
    %   the check and it works, mention it in the discussion/outlook as
    %   "re-analysis with the improved estimators may make the strap
    %   concept viable"; the methods chapter only needs the historical
    %   reason for the design change.
  \item Lessons carried forward: (i) compliant media work at metre scale but
    not at cm scale with velocity-derived strain; (ii) the reference must be
    a *measured boundary displacement*, not a derived field quantity —
    both directly motivate the final design.
\end{itemize}

\begin{figure}[htb]
  \centering
  \figplaceholder{(a) sweep space–time plot with migrating node (arrow);
    (b) velocity dispersion c(f) from node tracking (+ f–k slope)}
  \caption{Ratchet-strap characterisation via standing-wave nodes.
    [Condensed from your interim report.]}
  \label{fig:methods:strap}
\end{figure}

% ─────────────────────────────────────────────────────────────────────────────
\subsection{Main Experiment: Free-Hanging Cable Segments}
\label{sec:methods:campaign}

\subsubsection{Setup and Geometry}
\begin{itemize}
  \item [KEEP] your design paragraph: one cable end connected directly to
    the shaker source, the other clamped at a fixed distance → free-hanging
    segment strained under controlled conditions; a movable perpendicular
    beam screwed onto the profile sets the gap length; cable clamped at
    both shaker and fixed end.
  \item Gap lengths: 5, 10, 15\,cm; for each cable and gap a taut and an
    intended-sag configuration.
    % → GAP TO FILL (not in your notes): HOW the sag variants were produced
    %   (extra cable length fed in? mm of slack? how controlled/measured
    %   on-site?) and how reproducible that was. One short paragraph — the
    %   sag value itself is measured later from the scan geometry
    %   (processing chapter), so here only the physical procedure matters.
  \item Cables sprayed with chalk for reflectivity [KEEP].
  \item Coordinate offset correction (15.9 / 6 / 6 cm) — cross-ref the LDV
    subsection rather than repeating.
  \item Sketch + photo placeholders below [your notes ask for both].
\end{itemize}

\begin{figure}[htb]
  \centering
  \figplaceholder{Technical sketch: profile, shaker, connector, cable, gap
    L, clamps, measured geometric distances, local coordinate axes}
  \caption{Geometry of the free-hanging-cable setup. [Your sketch note.]}
  \label{fig:methods:maingeom}
\end{figure}

\begin{figure}[htb]
  \centering
  \figplaceholder{Photos: full setup + close-up of a mounted cable (taut and
    sag variant)}
  \caption{Free-hanging-cable experiment. [Your photo note.]}
  \label{fig:methods:mainphoto}
\end{figure}

\subsubsection{Measurement Grid and Reference}
\begin{itemize}
  \item [KEEP]: outermost point at each cable end selected in the LDV
    software; straight-line grid between them coarsened to 17 points on the
    cable (9 points for the 5\,cm gap).
  \item [KEEP]: per configuration, an additional recording on top of the
    shaker connector = the boundary-motion reference (the
    \texttt{\_SHAKER} datasets used as $\delta L$ input in the processing).
  \item Median stacking over 5 repetitions per point (cross-ref equipment
    section).
\end{itemize}

\subsubsection{Acquisition Parameters}
% Your note asked for exactly this table; values from the dataset naming
% convention / config.py — verify the LDV range column per dataset group.
\begin{table}[H]
  \centering
  \begin{tabularx}{\textwidth}{l X}
    \toprule
    \textbf{Parameter} & \textbf{Value} \\
    \midrule
    Sampling rate & 5\,kHz \\
    Samples per record & 60\,000 (12\,s; linearity: 190\,000 = 38\,s) \\
    Averages per point & 5 (median stack) \\
    Excitation & log sweep 1–500\,Hz, 10\,s + buffer
      (Table~\ref{tab:methods:signals}) \\
    Signal amplitude & 0.5\,V \\
    Amplifier setting & R2 \\
    LDV velocity range & 500 / 1250 / 5000\,mm\,s$^{-1}$ (dataset-dependent;
      thin cables needed larger ranges) \\
    Points on cable & 17 (gap 10/15\,cm) / 9 (gap 5\,cm) \\
    \bottomrule
  \end{tabularx}
  \caption[Acquisition parameters]{Acquisition parameters of the main
    dataset. Deviations per dataset are encoded in the file names
    (Appendix: naming convention).}
  \label{tab:methods:acqparams}
\end{table}

\subsubsection{Configuration Matrix and Dataset Inventory}
\begin{itemize}
  \item [KEEP] your dataset count paragraph — but re-check the arithmetic:
    $7 \times 3 \times 2 \times 2 = 84$, $+\,4 \times 2 = 8$, $+\,1 = 93$,
    not 98 as written. Recount against the file listing (note: not every
    combination exists — e.g.\ some Cable-7 sag variants replaced by TAPE
    configurations, one repeated recording at a different LDV range) and
    state the realised number.
  \item Deviations from the full matrix worth one sentence each: C7 15\,cm
    required taping (TAPE variants); repeated recordings at different LDV
    ranges; the cable-free direct shaker reference recording.
  \item Dataset naming convention: one appendix table decoding
    \texttt{Cable\{N\}\_\{gap\}cm[\_Sag][\_TAPE]\_avg5\_05V\_R2\_fs5kHz\_%
60000nt\_LogSweep\_1\_500Hz\_12\_5s\_\{range\}mms\_m2} (+ \texttt{\_SHAKER}
    suffix for the reference channel).
  \item Linearity-test subset: 4 configurations (C-IDs: verify against
    original numbering — code labels Cable5/6/7\_5cm and Cable7\_10cm\_Sag),
    mono-frequency ramp signal, 38\,s records.
\end{itemize}

% Optional compact matrix table (ticks = acquired):
\begin{table}[H]
  \centering
  \begin{tabularx}{\textwidth}{l *{6}{X}}
    \toprule
    & \multicolumn{2}{c}{5\,cm} & \multicolumn{2}{c}{10\,cm}
      & \multicolumn{2}{c}{15\,cm} \\
    \textbf{Cable} & taut & sag & taut & sag & taut & sag \\
    \midrule
    C1 \ldots C7 & \multicolumn{6}{c}{\textit{[tick marks / footnotes for
      TAPE, repeats, linearity subset]}} \\
    \bottomrule
  \end{tabularx}
  \caption[Configuration matrix]{Acquired configurations (✓), with
    footnotes for deviations.}
  \label{tab:methods:matrix}
\end{table}

% Closing sentence of the chapter: hand-off to the processing chapter —
% "All recordings were processed with the pipeline described in
%  Section \ref{sec:processing}."

% ─────────────────────────────────────────────────────────────────────────────
% OPEN ITEMS (delete as you resolve them)
% ─────────────────────────────────────────────────────────────────────────────
% [ ] Confirm x-offset: processing used 15.9 cm (io.py) — re-measure /
%     document why the tape measure said 14.9.
% [ ] C5 linear weight: 6.67 vs computed 6.35 g/m — recompute density.
% [ ] Dataset count: formula gives 93, notes say 98 — recount files.
% [ ] Describe HOW intended sag was induced (procedure + control).
% [ ] Strap experiment: exact spatial sampling value; C1–C4 = which
%     numbering; acquisition table for appendix.
% [ ] Citation for MATRIX/CIWE platform + Polytec LDV (50 cm standoff).
% [ ] Cable construction details per cable (spec sheets) for
%     Table~\ref{tab:methods:cabletypes}.
% [ ] Verify E means/stds against lab sheets (values above from config.py,
%     mapped old→new numbering).
% [ ] Decide appendix set: weighing sheets, tensile curves, strap report,
%     naming convention, acquisition tables of preliminary experiments.
